\chapter{Course templates}
\label{ch:templates}

A template stores the \emph{layout} of a course: which rooms exist, in which order, with
which names, topics and settings, and how many group rooms. It contains no people and no
messages. Use templates to set up the same kind of course again, or to share a layout.

\section{Save a course as a template}

\begin{you}
!template export CS101 lecture-with-labs
\end{you}

The bot saves the layout of \cmd{CS101} as your template \cmd{lecture-with-labs} (without a
name: the course code) and sends it to you as a \file{.yaml} file. To get a saved template
again later, even if its course no longer exists:
\cmd{!template export --saved lecture-with-labs}.

\section{Use a template}

\begin{you}
!course create CS202 "Algorithms" --template lecture-with-labs
\end{you}

\cmd{!template list} shows your own templates (\emph{personal}) and the ones bot admins made
available to everyone (\emph{global}). \cmd{!template delete <name>} deletes one of yours;
courses created from it are not affected.

\section{Edit a template}

Download it with \cmd{!template export}, change the file in a text editor, and load it again:

\begin{you}
!template import lecture-with-labs
\end{you}

Then send the \file{.yaml} file. The bot checks it and, if something is wrong, says what and
where (e.g.\ \enquote{rooms[2]: unknown keys: titel}); send the corrected file without
repeating the command. Importing under an existing name replaces that template.
\cmd{--global} makes a template available to all professors (bot admins only).

\subsection*{The format}

\begin{codefile}[lecture-with-labs.yaml]
version: 1
space: {name: "{code} {title} ({semester})", topic: "Lecture with weekly labs"}
defaults: {join_rule: restricted, encryption: true, history_visibility: shared}
rooms:
  - {name: "Announcements", power_levels: {events_default: 50}}
  - {name: "General"}
  - {name: "Questions"}
  - {name: "Lab sessions", topic: "Questions about the weekly lab"}
  - {name: "Staff", members: staff}
groups: {count: 0, name: "Group {n}"}
\end{codefile}

\begin{center}
\begin{tabularx}{\textwidth}{@{}lL@{}}
\toprule
\textbf{Key} & \textbf{Meaning} \\
\midrule
\cmd{version} & always \cmd{1} \\
\cmd{space.name} & name of the space; \cmd{{code}}, \cmd{{title}} and \cmd{{semester}} are filled in \\
\cmd{space.topic} & optional description of the space \\
\cmd{defaults} & settings for every room unless the room says otherwise (below) \\
\cmd{rooms} & the rooms in order (1 to 50); each needs a \cmd{name} \\
\quad\cmd{topic} & optional room description \\
\quad\cmd{members} & \cmd{all} (default) or \cmd{staff}: only owner and TAs, invite-only \\
\quad\cmd{join_rule} & \cmd{restricted}: everyone in the course space can join (default);
                       \cmd{invite}: by invitation only; \cmd{public}: anyone on the server \\
\quad\cmd{encryption} & \cmd{true} (default) or \cmd{false} \\
\quad\cmd{history_visibility} & \cmd{shared}: new members see older messages (default);
                                \cmd{invited} or \cmd{joined}: only from that point on \\
\quad\cmd{power_levels} & deviations from the standard permissions, e.g.\
                          \cmd{events_default: 50}: only staff can write \\
\cmd{groups} & \cmd{count}: group rooms created with the course (0 to 50);
               \cmd{name}: their name, must contain \cmd{{n}} \\
\bottomrule
\end{tabularx}
\end{center}

Allowed \cmd{power_levels} keys are \cmd{users_default}, \cmd{events_default},
\cmd{state_default}, \cmd{invite}, \cmd{kick}, \cmd{ban} and \cmd{redact}, each 0--100.
Students have level 0, TAs 50 and the owner 100.
